% Revised paragraph - more careful about what we claim vs. speculate
% Addresses composite objects without overreaching

An important question arises for composite objects: when an asteroid contains $10^{33}$ atoms, how does the unified gravitational envelope maintain $k=4$ closure? The field equations ensure that the field is \emph{not} a simple average over the mass distribution—this would wash out the topological structure encoded in the tetrad bases. The field \emph{is} isomorphic to the object's configuration: it couples to the actual mass distribution through the energy-momentum tensor $T_{\mu\nu}$ in the field equations. What remains an open question is the internal mechanism: does the field maintain $k=4$ topology solely at the outer envelope boundary, or does it require additional structure at internal scale-transition surfaces? One possibility is that field closure operates at radial domain boundaries—transition regions where the characteristic scale changes—with $k=4$ tetrahedral structure enforced at each such boundary. This would create a hierarchical mesh where local topology at each scale couples to the envelope-level topology. However, the field equations do not uniquely determine this internal structure; they constrain only that $\det(e) \neq 0$ must hold everywhere, which requires $k=4$ basis vectors at every point, and that the envelope-level field exhibits $k=4$ topology when coupling to external fields. The physical object may be elongated and irregular, yet the envelope field—the effective gravitational potential seen by external observers—must satisfy the $k=4$ topological constraint. Whether this constraint propagates through explicit radial boundaries or through other topological mechanisms is beyond the scope of this analysis, but the key result is that $k=4$ at the envelope level ensures the Poisson equation (\ref{eq:poisson-numeric}) applies universally.
